% =============================================================================
%  lec16.tex —— 第 16 讲 t-结构（底本 PDF 第 124–130 页）
% =============================================================================
\lecture{16}{t-结构}
\begin{center}\small\color{black!55}（底本 PDF 第 124–130 页）\end{center}

本讲介绍稳定 $\infty$-范畴上的 t-结构。

\section{定义}

\nitem\label{L16-1.1} 设 $\cC$ 是稳定 $\infty$-范畴。$\cC$ 上的 t-结构就是三角范畴 $\mathrm{h}\cC$ 上的 t-结构。

\begin{definition}\label{L16-1.2}
设 $\cC$ 是稳定 $\infty$-范畴。$\cC$ 上的一个 t-结构定义为一对 $\cC$ 的全子 $\infty$-范畴 $(\cC^{\le0},\cC^{\ge0})$，使得当我们记
\[
\cC^{\le n}:=\cC^{\le0}[-n],\qquad \cC^{\ge n}:=\cC^{\ge0}[-n]
\]
时，有：
\begin{enumerate}[label=(\arabic*),start=0,leftmargin=2.4em,itemsep=0.1em]
\item $\cC^{\le0}$ 与 $\cC^{\ge0}$ 都在 $\cC$ 中的同构下稳定。
\item 对 $X\in\cC^{\le0}$ 与 $Y\in\cC^{\ge1}$，有 $\operatorname{Maps}_{\cC}(X,Y)\simeq\{0\}$。
\item 有包含 $\cC^{\le-1}\subseteq\cC^{\le0}$ 与 $\cC^{\ge1}\subseteq\cC^{\ge0}$。
\item 对任意 $X\in\cC$，存在纤维—余纤维序列 $X'\to X\to X''$，使得 $X'\in\cC^{\le0}$ 且 $X''\in\cC^{\ge1}$。
\end{enumerate}
\end{definition}

\begin{warning}\label{L16-1.3}
我们使用上同调约定。为与同调约定下的记号比较，令
\[
\cC_{\le n}:=\cC^{\ge-n},\qquad \cC_{\ge n}:=\cC^{\le-n}.
\]
\end{warning}

\begin{exercise}\label{L16-1.4}
证明对应
\[
(\cC^{\le0},\cC^{\ge0})\mapsto(\mathrm{h}\cC^{\le0},\mathrm{h}\cC^{\ge0})
\]
给出 $\cC$ 上的 t-结构与三角范畴 $\mathrm{h}\cC$ 上的 t-结构之间的双射。
\end{exercise}

\begin{exercise}\label{L16-1.5}
证明对应
\[
(\cC^{\le0},\cC^{\ge0})\mapsto((\cC^{\ge0})^{\op},(\cC^{\le0})^{\op})
\]
给出 $\cC$ 上的 t-结构与 $\cC^{\op}$ 上的 t-结构之间的双射。
\end{exercise}

\begin{exercise}\label{L16-1.6}
证明：对 $m<n$，有 $\cC^{\le m}\cap\cC^{\ge n}\simeq\{0\}$。
\end{exercise}

\begin{lemma}\label{L16-1.7}
设 $n\ge0$。对任意 $X\in\cC^{\le0}$ 与 $Y\in\cC^{\ge-n}$，映射空间 $\operatorname{Maps}_{\cC}(X,Y)$ 是同伦 $n$-型。
\end{lemma}

\begin{proof}[证明梗概]
首先注意 $\operatorname{Maps}_{\cC}(X,Y)$ 的连通分支彼此弱同伦等价，因为它是 $\operatorname{Maps}_{\cC}(X,\Sigma Y)$ 的环路空间。因此只需证明 $\pi_m\operatorname{Maps}_{\cC}(X,Y)\simeq0$（$m>n$），基点取 $0\in\operatorname{Maps}_{\cC}(X,Y)$。这可由下述观察得到：$\Omega^{n+1}\operatorname{Maps}_{\cC}(X,Y)\simeq\operatorname{Maps}_{\cC}(X,\Omega^{n+1}Y)$ 弱可缩，因为 $\Omega^{n+1}Y\in\cC^{\ge1}$。
\end{proof}

\begin{definition}\label{L16-1.8}
设 $\cC$ 是带 t-结构的稳定 $\infty$-范畴。该 t-结构的\textbf{心}定义为
\[
\cC^{\heartsuit}:=\cC^{\le0}\cap\cC^{\ge0}.
\]
\end{definition}

\begin{theorem}\label{L16-1.9}
$\infty$-范畴 $\cC^{\heartsuit}$ 是通常的 Abel 范畴。此外，$0\to X''\to X\to X'\to 0$ 是 $\cC^{\heartsuit}$ 中的短正合序列，当且仅当 $X''\to X\to X'$ 是 $\cC$ 的、含于 $\cC^{\heartsuit}$ 中的纤维—余纤维序列。
\end{theorem}

\begin{proof}[证明梗概]
由引理~\ref{L16-1.7}，$\cC^{\heartsuit}$ 是通常范畴。由此可将它与 $\mathrm{h}\cC^{\le0}\cap\mathrm{h}\cC^{\ge0}$ 等同，而后者是三角范畴 $\mathrm{h}\cC$ 的心。于是所述断言由三角范畴中相应的熟知断言推出。
\end{proof}

\begin{warning}\label{L16-1.10}
t-结构不由其心决定。稳定 $\infty$-范畴中的大部分信息都不能从心中恢复。例如，对 $X,Y\in\cC^{\heartsuit}$，群 $\operatorname{Ext}^{\bullet}_{\cC}(X,Y)$ 与 $\operatorname{Ext}^{\bullet}_{\cC^{\heartsuit}}(X,Y)$ 一般不同构。
\end{warning}

\begin{definition}\label{L16-1.11}
设 $\cC$ 与 $\cD$ 是带 t-结构的稳定 $\infty$-范畴。对正合函子 $F\colon\cC\to\cD$，我们说
\begin{itemize}[leftmargin=2.2em,itemsep=0.1em]
\item 函子 $F$ 是\textbf{左 t-正合}的，若它把 $\cC^{\ge0}$ 映入 $\cD^{\ge0}$；
\item 函子 $F$ 是\textbf{右 t-正合}的，若它把 $\cC^{\le0}$ 映入 $\cD^{\le0}$；
\item 函子 $F$ 是\textbf{t-正合}的，若它既左 t-正合又右 t-正合。
\end{itemize}
\end{definition}

\begin{exercise}\label{L16-1.12}
证明：左 t-正合函子的左伴随是右 t-正合的。
\end{exercise}

\begin{exercise}\label{L16-1.13}
构造一个保持心、但不 t-正合的正合函子 $F\colon\cC\to\cD$。
\end{exercise}

\section{上同调}

\nitem\label{L16-2.1}\textbf{命题—定义}\quad 设 $\cC$ 是带 t-结构的稳定 $\infty$-范畴。对任意 $n$，包含函子 $\cC^{\le n}\to\cC$ 有右伴随 $\tau^{\le n}$，包含函子 $\cC^{\ge n}\to\cC$ 有左伴随 $\tau^{\ge n}$。这些函子称为该 t-结构的\textbf{截断函子}。

\begin{remark}\label{L16-2.2}
上述左右可从下述图记忆：$\cC^{\le0}\rightleftarrows\cC\rightleftarrows\cC^{\ge1}$。
\end{remark}

\begin{exercise}\label{L16-2.3}
通过验证下述断言来证明上述命题。设 $X\in\cC$，$X'\to X\to X''$ 是任意纤维—余纤维序列，其中 $X'\in\cC^{\le0}$ 且 $X''\in\cC^{\ge1}$，则
\begin{enumerate}[label=(\arabic*),leftmargin=2.4em,itemsep=0.1em]
\item 对任意 $Y\in\cC^{\le0}$，态射 $X'\to X$ 诱导等价
\[
\operatorname{Maps}_{\cC}(Y,X')\xrightarrow{\ \simeq\ }\operatorname{Maps}_{\cC}(Y,X).
\]
\item 对任意 $Z\in\cC^{\ge1}$，态射 $X\to X''$ 诱导等价
\[
\operatorname{Maps}_{\cC}(X,Z)\xrightarrow{\ \simeq\ }\operatorname{Maps}_{\cC}(X'',Z).
\]
\end{enumerate}
并推出 $X'\simeq\tau^{\le0}X$ 与 $X''\simeq\tau^{\ge1}X$。
\end{exercise}

\begin{exercise}\label{L16-2.4}
我们有笛卡尔方块
\begin{center}
\begin{tikzcd}[column sep=3.2em,row sep=2.6em]
\cC^{\le0} \arrow[r, hook, "\subseteq"] \arrow[d] & \cC \arrow[d, "\tau^{\ge1}"] \\
\{0\} \arrow[r, hook, "\subseteq"] & \cC^{\ge1}
\end{tikzcd}
\qquad
\begin{tikzcd}[column sep=3.2em,row sep=2.6em]
\cC^{\ge0} \arrow[r, hook, "\subseteq"] \arrow[d] & \cC \arrow[d, "\tau^{\le-1}"] \\
\{0\} \arrow[r, hook, "\subseteq"] & \cC^{\le-1}
\end{tikzcd}
\end{center}
特别地，$\cC^{\le0}$ 与 $\cC^{\ge0}$ 互相决定。
\end{exercise}

\begin{exercise}\label{L16-2.5}
全子 $\infty$-范畴 $\cC^{\le0}\subseteq\cC$ 在余极限下稳定，而 $\cC^{\le0}\subseteq\cC$ 在极限下稳定。
\end{exercise}

\yizhu{原文此处两次写作 $\cC^{\le0}\subseteq\cC$：前一处说它在余极限下稳定，后一处说它在极限下稳定。由命题—定义~\ref{L16-2.1}，截断函子的左右伴随性表明在极限下稳定的应是 $\cC^{\ge0}\subseteq\cC$，第二处疑为笔误；译文照录原文。\textbf{[来源：I]}}

\begin{exercise}\label{L16-2.6}
$\cC^{\le0}$ 中的极限同构于 $\cC$ 中相应极限的 $\tau^{\le0}$-截断。对偶地，$\cC^{\ge0}$ 中的余极限同构于 $\cC$ 中相应余极限的 $\tau^{\ge0}$-截断。
\end{exercise}

\nitem\label{L16-2.7} 下述练习说明：对纤维—余纤维序列 $X'\to X\to X''$，两项的振幅估计给出第三项的估计。

\begin{exercise}\label{L16-2.8}
设 $X'\to X\to X''$ 是纤维—余纤维序列。
\begin{enumerate}[label=(\arabic*),leftmargin=2.4em,itemsep=0.1em]
\item 若 $X',X''\in\cC^{\le0}$，则 $X\in\cC^{\le0}$。
\item 若 $X\in\cC^{\le0}$ 且 $X''\in\cC^{\le-1}$，则 $X'\in\cC^{\le0}$。
\item 若 $X\in\cC^{\le0}$ 且 $X'\in\cC^{\le1}$，则 $X''\in\cC^{\le0}$。
\end{enumerate}
特别地，全子 $\infty$-范畴 $\cC^{\le0}\subseteq\cC\supseteq\cC^{\ge0}$ 在扩张下稳定。
\end{exercise}

\begin{exercise}\label{L16-2.9}
举例说明上述估计是最优的。
\end{exercise}

\begin{exercise}\label{L16-2.10}
截断函子 $\tau_{\le\bullet}$ 与 $\tau_{\ge\bullet}$ 彼此交换。
\end{exercise}

\begin{remark}\label{L16-2.11}
上述练习的确切含义如下。对 $m$ 与 $n$，有
\begin{itemize}[leftmargin=2.2em,itemsep=0.1em]
\item $\tau_{\ge m}\circ\tau_{\ge n}\simeq\tau_{\ge\max\{m,n\}}$，$\tau_{\le m}\circ\tau_{\le n}\simeq\tau_{\le\max\{m,n\}}$。
\item 交换方块
\begin{center}
\begin{tikzcd}[column sep=4.0em,row sep=2.6em]
\cC^{\le m}\cap\cC^{\ge n} \arrow[r, hook, "\subseteq"] \arrow[d, hook, "\subseteq"'] & \cC^{\le m} \arrow[d, hook, "\subseteq"] \\
\cC^{\ge n} \arrow[r, hook, "\subseteq"] & \cC
\end{tikzcd}
\end{center}
沿水平方向左可伴随，且诱导的交换方块
\begin{center}
\begin{tikzcd}[column sep=4.0em,row sep=2.6em]
\cC^{\le m}\cap\cC^{\ge n} \arrow[d, hook, "\subseteq"'] & \cC^{\le m} \arrow[l, "\tau^{\ge n}"'] \arrow[d, hook, "\subseteq"] \\
\cC^{\ge n} & \cC \arrow[l, "\tau^{\ge n}"']
\end{tikzcd}
\end{center}
沿竖直方向右可伴随，即诱导交换方块
\begin{center}
\begin{tikzcd}[column sep=4.0em,row sep=2.6em]
\cC^{\le m}\cap\cC^{\ge n} & \cC^{\le m} \arrow[l, "\tau^{\ge n}"'] \\
\cC^{\ge n} \arrow[u, "\tau^{\le m}"] & \cC \arrow[l, "\tau^{\ge n}"'] \arrow[u, "\tau^{\le m}"]
\end{tikzcd}
\end{center}
\end{itemize}
\end{remark}

\yizhu{注~\ref{L16-2.11} 第一个要点的第二个等式，右端原文亦作 $\max\{m,n\}$；但 $\tau_{\le m}\circ\tau_{\le n}$ 等同于 $\tau_{\le\min\{m,n\}}$，故此处应为 $\min$。译文照录原文。\textbf{[来源：B]}}

\begin{definition}\label{L16-2.12}
设 $\cC$ 是带 t-结构的稳定 $\infty$-范畴。考虑函子
\[
H^n\colon \cC\xrightarrow{\ \tau^{\le n}\circ\tau^{\ge n}\ }\cC^{\le n}\cap\cC^{\ge n}\xrightarrow{\ [-n]\ }\cC^{\heartsuit}.
\]
对 $X\in\cC$，我们称 $H^n(X)\in\cC^{\heartsuit}$ 为 $X$ 的\textbf{第 $n$ 个上同调}。
\end{definition}

\begin{warning}\label{L16-2.13}
可能出现 $H^n(X)\simeq0$ 而 $X$ 不同构于 $0$ 的情况。例如，$\cC^{\le0}:=\cC$ 与 $\cC^{\ge0}:=\{0\}$ 给出一个 t-结构，其心 $\cC^{\heartsuit}\simeq\{0\}$。在下述情形中有大量有趣的例子：
\begin{itemize}[leftmargin=2.2em,itemsep=0.1em]
\item 非正则代数几何；
\item 无限型代数几何；
\item 无限维表示论；
\item ……
\end{itemize}
\end{warning}

\nitem\label{L16-2.14} 下述结果可由三角范畴中相应的熟知结果推出。

\begin{proposition}\label{L16-2.15}
设 $\cC$ 是带 t-结构的稳定 $\infty$-范畴。对任意纤维—余纤维序列 $X'\to X\to X''$，我们在 $\cC^{\heartsuit}$ 中有长正合序列
\[
\cdots\to H^n(X')\to H^n(X)\to H^n(X'')\xrightarrow{\ \delta\ }H^{n+1}(X')\to\cdots,
\]
其中连接态射 $\delta$ 由态射 $X''\to X'[1]$ 诱导。
\end{proposition}

\nitem\label{L16-2.16} 下述结果把 t-正合函子与心之间的正合函子联系起来。

\begin{exercise}\label{L16-2.17}
设 $F\colon\cC\to\cD$ 是左 t-正合函子。证明复合
\[
H^0F\colon \cC^{\heartsuit}\to\cC^{\ge0}\xrightarrow{\ F\ }\cD^{\ge0}\xrightarrow{\ \tau^{\le0}\ }\cD^{\heartsuit}
\]
是左正合的。
\end{exercise}

\begin{warning}\label{L16-2.18}
与警告~\ref{L16-1.10} 一样，即便是 t-正合函子 $F$，也不能从 $H^0F$ 恢复。
\end{warning}

\begin{remark}\label{L16-2.19}
下次我们将定义 Abel 范畴 $\cA$ 的各种版本的导出 $\infty$-范畴，它们带有 t-结构，其心可等同于 $\cA$。按定义，关于这些导出 $\infty$-范畴的全部信息都可从 $\cA$ 恢复。

在某些假设下，从这些导出 $\infty$-范畴出发的左/右 t-正合函子 $F$ 可等同于 $H^0F$ 的右/左导出函子。
\end{remark}

\section{有界、分离与完备}

\begin{definition}\label{L16-3.1}
设 $\cC$ 是带 t-结构的稳定 $\infty$-范畴。对对象 $X\in\cC$，
\begin{itemize}[leftmargin=2.2em,itemsep=0.1em]
\item 称 $X$ 是\textbf{连通}的，若 $X\in\cC^{\le0}$；
\item 称 $X$ 是\textbf{余连通}的，若 $X\in\cC^{\ge0}$；
\item 称 $X$ 是\textbf{$n$-连通}的，若 $X\in\cC^{\le-n}$；
\item 称 $X$ 是\textbf{$n$-余连通}的，若 $X\in\cC^{\ge-n}$；
\item 称 $X$ 是\textbf{最终连通}的，或\textbf{右有界}的，若 $X\in\cC^{-}:=\bigcup_n\cC^{\le n}$；
\item 称 $X$ 是\textbf{最终余连通}的，或\textbf{左有界}的，若 $X\in\cC^{+}:=\bigcup_n\cC^{\ge n}$；
\item 称 $X$ 是\textbf{有界}的，若 $X$ 既左有界又右有界：$X\in\cC^{b}:=\cC^{+}\cap\cC^{-}$。
\end{itemize}
\end{definition}

\begin{warning}\label{L16-3.2}
与警告~\ref{L16-2.13} 一样，上述性质不能通过上同调检验。
\end{warning}

\begin{exercise}\label{L16-3.3}
证明：若 $X$ 左有界，则 $X\in\cC^{\ge0}$ 当且仅当对 $i<0$ 有 $H^i(X)\simeq0$。
\end{exercise}

\begin{definition}\label{L16-3.4}
设 $\cC$ 是带 t-结构的稳定 $\infty$-范畴。
\begin{itemize}[leftmargin=2.2em,itemsep=0.1em]
\item 称 $\cC$ 是\textbf{右有界}的，若 $\cC\simeq\cC^{-}$。
\item 称 $\cC$ 是\textbf{左有界}的，若 $\cC\simeq\cC^{+}$。
\item 称 $\cC$ 是\textbf{有界}的，若 $\cC\simeq\cC^{b}$。
\item 称 $\cC$ 是\textbf{右分离}的，若
\[
\cC^{\ge\infty}:=\bigcap_n\cC^{\ge n}\simeq\{0\}.
\]
\item 称 $\cC$ 是\textbf{左分离}的，若
\[
\cC^{\le-\infty}:=\bigcap_n\cC^{\le n}\simeq\{0\}.
\]
\end{itemize}
\end{definition}

\begin{definition}\label{L16-3.5}
设 $\cC$ 是带 t-结构的稳定 $\infty$-范畴。$\cC$ 的\textbf{左完备化} $\widehat{\cC}$ 定义为下述图的极限
\[
\cdots\xrightarrow{\tau}\cC^{\ge-2}\xrightarrow{\tau}\cC^{\ge-1}\xrightarrow{\tau}\cC^{\ge0}\xrightarrow{\tau}\cdots.
\]
称 $\cC$ 是\textbf{左完备}的，若函子 $\cC\to\widehat{\cC}$ 是等价。

对偶地，\textbf{右完备化}定义为
\[
\cdots\xrightarrow{\tau}\cC^{\le2}\xrightarrow{\tau}\cC^{\le1}\xrightarrow{\tau}\cC^{\le0}\xrightarrow{\tau}\cdots
\]
的极限。
\end{definition}

\begin{exercise}\label{L16-3.6}
证明 $\cC$ 的右完备化等价于 $\mathrm{Sptr}(\cC^{\le0})$。
\end{exercise}

\begin{exercise}\label{L16-3.7}
证明左完备化 $\widehat{\cC}$ 是稳定的，且函子 $\cC\to\widehat{\cC}$ 正合。
\end{exercise}

\begin{exercise}\label{L16-3.8}
证明左完备化 $\widehat{\cC}$ 具有本质唯一的 t-结构，使得 $\cC^{\ge0}\xrightarrow{\ \simeq\ }(\widehat{\cC})^{\ge0}$。
\end{exercise}

\begin{exercise}\label{L16-3.9}
证明 $\cC$ 左分离，当且仅当函子 $\cC\to\widehat{\cC}$ 检测零对象。特别地，若 $\cC$ 左完备，则它左分离。
\end{exercise}

\begin{warning}\label{L16-3.10}
左分离的 t-结构未必左完备。例如，任意左有界的 t-结构都是左分离的，但几乎从不左完备。
\end{warning}

\begin{exercise}\label{L16-3.11}
设 $\cC$ 具有可数积，且 $\cC^{\le0}\subseteq\cC$ 在可数积下稳定\footnote{这些条件常被称作：取可数积是右 t-正合的。}。则 $\cC$ 左分离，当且仅当它左完备。
\end{exercise}

\begin{exercise}\label{L16-3.12}
证明 $\cC$ 左完备，当且仅当它是 Postnikov 完备的。后者意即
\begin{itemize}[leftmargin=2.2em,itemsep=0.1em]
\item 每个对象 $X\in\cC$ 都是其 Postnikov 塔的极限：
\[
X\simeq\lim[\cdots\to\tau^{\ge n-1}X\to\tau^{\ge n}X\to\tau^{\ge n+1}X\to\cdots]
\]
\item $\cC$ 中任意 Postnikov 塔都收敛。
\end{itemize}
换言之，任意一族的 $X_n\in\cC^{\ge n}$ 连同同构 $\tau^{\ge n}X_{n-1}\xrightarrow{\ \simeq\ }X_n$ 都是
\[
X:=\lim[\cdots\to X_{n-1}\to X_n\to X_{n+1}\to\cdots]\in\cC
\]
的 Postnikov 塔。
\end{exercise}

\section{可表现稳定 $\infty$-范畴上的 t-结构}

\nitem\label{L16-4.1}\textbf{命题—定义}（HA.1.4.4.13）\quad 设 $\cC$ 是可表现稳定 $\infty$-范畴，带 t-结构。下述条件等价：
\begin{itemize}[leftmargin=2.2em,itemsep=0.1em]
\item $\infty$-范畴 $\cC^{\le0}$ 可表现。
\item $\infty$-范畴 $\cC^{\le0}$ 可及。
\item $\infty$-范畴 $\cC^{\ge0}$ 可表现。
\item $\infty$-范畴 $\cC^{\ge0}$ 可及。
\item 复合 $\cC\xrightarrow{\tau^{\le0}}\cC^{\le0}\to\cC$ 可及。
\item 复合 $\cC\xrightarrow{\tau^{\ge0}}\cC^{\ge0}\to\cC$ 可及。
\end{itemize}
称 $\cC$ 上的 t-结构是\textbf{可及的}，若它满足上述条件。

\begin{proposition}\label{L16-4.2}（HA.1.4.4.11）
设 $\cC$ 是可表现稳定 $\infty$-范畴。全子 $\infty$-范畴 $\cC'\subseteq\cC$ 决定一个可及 t-结构且 $\cC^{\le0}\simeq\cC'$，当且仅当
\begin{itemize}[leftmargin=2.2em,itemsep=0.1em]
\item $\cC'$ 可及；
\item $\cC'\subseteq\cC$ 在余极限与扩张下封闭。
\end{itemize}
\end{proposition}

\begin{remark}\label{L16-4.3}
实践中，可对由一小族对象在小余极限与扩张下生成的最小子 $\infty$-范畴 $\cC'$ 应用上述命题。这样的 $\cC'$ 总是可及的。
\end{remark}

\begin{definition}\label{L16-4.4}
设 $\cC$ 是可表现稳定 $\infty$-范畴。称 $\cC$ 上的 t-结构与\textbf{滤过余极限相容}，若取滤过余极限是 t-正合的。
\end{definition}

\begin{exercise}\label{L16-4.5}
设 $\cC$ 是带 t-结构的可表现稳定 $\infty$-范畴。证明下述条件等价：
\begin{itemize}[leftmargin=2.2em,itemsep=0.1em]
\item 该 t-结构与滤过余极限相容。
\item 嵌入 $\cC^{\ge0}\subseteq\cC$ 在滤过余极限下稳定。
\item $\tau^{\le0}$ 保持滤过余极限。
\end{itemize}
\end{exercise}

\begin{remark}\label{L16-4.6}
注意 $\tau^{\ge0}$ 与 $\cC^{\le0}\subseteq\cC$ 总保持任意余极限，因为它们有右伴随。
\end{remark}

\begin{definition}\label{L16-4.7}
设 $\cC$ 是可表现稳定 $\infty$-范畴。称 $\cC$ 上的 t-结构是\textbf{紧生成}的，若 $\cC^{\le0}$ 紧生成，且函子 $\cC^{\le0}\to\cC$ 保持紧对象。
\end{definition}

\begin{exercise}\label{L16-4.8}
证明任意紧生成的 t-结构可及且与滤过余极限相容。
\end{exercise}

\begin{exercise}\label{L16-4.9}
设 $\cC$ 是带右完备 t-结构的可表现稳定 $\infty$-范畴。证明下述条件等价：
\begin{itemize}[leftmargin=2.2em,itemsep=0.1em]
\item 该 t-结构紧生成。
\item $\infty$-范畴 $\cC^{\le0}$ 紧生成。
\end{itemize}
此外，在这些条件下，$\cC$ 紧生成，且 $\cC$ 中的紧对象恰为 $Y[n]$，其中 $Y$ 是 $\cC^{\le0}$ 中的紧对象、$n$ 为整数。
\end{exercise}

\begin{exercise}\label{L16-4.10}
设 $\cC_0$ 是带 t-结构的小稳定 $\infty$-范畴，且该 t-结构有界。证明 $\cC:=\Ind(\cC_0)$ 具有紧生成的右完备 t-结构，且 $\cC^{\le0}:=\Ind(\cC_0^{\le0})$。
\end{exercise}

\begin{warning}\label{L16-4.11}
上述 $\Ind(\cC_0)$ 上的 t-结构一般不是左完备的。在警告~\ref{L16-2.13} 所列的情形中有大量有趣的例子。
\end{warning}

\begin{warning}\label{L16-4.12}
$\cC$ 中的紧对象未必在截断下封闭。因此并非所有紧生成的 t-结构都来自 ind-完备化。在警告~\ref{L16-2.13} 所列的情形中有大量有趣的例子。
\end{warning}

\section{$\mathrm{Sptr}$ 上的 t-结构}

\nitem\label{L16-5.1}\textbf{命题—构造}（HA.1.4.3.4）\quad 设 $\cC$ 是带基点的可表现 $\infty$-范畴。我们有以下述方式给出的 $\mathrm{Sptr}(\cC)$ 上的可及 t-结构。
\begin{itemize}[leftmargin=2.2em,itemsep=0.1em]
\item 令 $\mathrm{Sptr}(\cC)^{\ge1}$ 为满足 $\Omega^{\infty}(X)\simeq0$ 的对象 $X$ 所成的全子 $\infty$-范畴。
\item 令 $\mathrm{Sptr}(\cC)^{\le0}$ 为由函子 $\Sigma^{\infty}$ 的本质像在扩张与小余极限下生成的全子 $\infty$-范畴。
\end{itemize}

\begin{proposition}\label{L16-5.2}（HA.1.4.3.6）
设 $\mathrm{Sptr}$ 带有上述 t-结构。则
\begin{enumerate}[label=(\arabic*),leftmargin=2.4em,itemsep=0.1em]
\item 该 t-结构紧生成\footnote{紧生成性在 HA 中没有证明，但它可由 $\Spc_*$ 紧生成这一事实推出。}、左完备且右完备。
\item Eilenberg–MacLane 谱函子 $A\mapsto HA$ 给出等价
\[
\Ab\xrightarrow{\ \simeq\ }\mathrm{Sptr}^{\heartsuit}.
\]
\end{enumerate}
\end{proposition}

\begin{exercise}\label{L16-5.3}
证明上同调函子 $H^i\colon\mathrm{Sptr}\to\mathrm{Sptr}^{\heartsuit}$ 可通过等价 $\Ab\xrightarrow{\ \simeq\ }\mathrm{Sptr}^{\heartsuit}$ 与 $\pi_{-i}$ 等同。
\end{exercise}

\begin{warning}\label{L16-5.4}
设 $A_1$ 与 $A_2$ 是 Abel 群。一般地，扩张群 $\operatorname{Ext}^i_{\mathrm{Sptr}}(HA_1,HA_2)$ 与 $\operatorname{Ext}^i(A_1,A_2)$ 不同构。例如，分次环 $\operatorname{Ext}^{\bullet}_{\mathrm{Sptr}}(HF_q,HF_q)$ 是模 $p$ 的 Steenrod 代数。
\end{warning}

\nitem\label{L16-5.5} 下次我们将构造自然同态
\[
\operatorname{Ext}^i(A_1,A_2)\to\operatorname{Ext}^i_{\mathrm{Sptr}}(HA_1,HA_2),
\]
它应被视为沿 $\Spec\Z\to\Spec\mathbb{S}$ 的导出正像。

\appendixsection{预稳定 $\infty$-范畴}

\nitem\label{L16-A.1} 设 $\cC$ 是带 t-结构的稳定 $\infty$-范畴。有时研究 $\cC^{\le0}$ 比研究 $\cC$ 更为方便。可以找到若干组公理来刻画形如 $\cC^{\le0}$ 的 $\infty$-范畴。这样的 $\infty$-范畴称为\textbf{预稳定 $\infty$-范畴}。

\nitem\label{L16-A.2} \textbf{建议阅读}\quad SAG.C。

\yizhu{“SAG”指 Lurie 的《Spectral Algebraic Geometry》（谱代数几何），“C”为其章节编号；原文此行仅给出这一缩写作为建议阅读，未展开说明。\textbf{[来源：B]}}
